\section{傅里叶级数}

\subsection{三角函数空间}

某个区间上的函数可以构成一个 Euclid 空间。其中加法和数乘就是函数的运算，而内积定义为：

\begin{definition}[函数的内积]
    对于 $[l, r]$ 上的两个函数 $f(x), g(x)$，我们定义它们的内积为
    $$
    \ab\langle f, g \rangle = \int_l^r f(x) g(x) \dd{x}
    $$
\end{definition}

有了内积之后，就可以定义正交：

\begin{definition}[正交]
    若 $\ab\langle f, g \rangle = 0$，则称 $f$ 与 $g$ 正交。
\end{definition}

三角函数系是指下面的函数线性组合得到的函数族：
$$
1,\ \sin n x,\ \cos n x,\ (n \in \mathbb{N}^+)
$$

在三角函数系中，我们定义函数的正交取区间 $[-\pi, \pi]$。那么我们可以发现：

\begin{itemize}
    \item
    $$
    \int_{-\pi}^{\pi} \sin n x \dd{x} = \int_{-\pi}^{\pi} \cos n x \dd{x} = 0
    $$
    \item
    $$
    \int_{-\pi}^{\pi} \sin m x \sin n x \dd{x} = \pi [m = n]
    $$
    \item
    $$
    \int_{-\pi}^{\pi} \cos m x \cos n x \dd{x} = \pi [m = n]
    $$
    \item
    $$
    \int_{-\pi}^{\pi} \sin m x \cos n x \dd{x} = 0
    $$
\end{itemize}

\subsection{Fourier 级数}

\begin{definition}
    一般地，\textbf{三角级数}定义为如下形式的级数：
    $$
    f(x) = \frac{a_0}{2} + \sum_{n=1}^{\infty} (a_n \cos n x + b_n \sin n x)
    $$
\end{definition}

根据函数项级数的 Weierstrass 定理，我们可以知道，若级数
$$
\sum_{n = 1}^{\infty} (\abs{a_n} + \abs{b_n})
$$
收敛，那么对应的三角级数在 $[-\pi, \pi]$ 上一致收敛且绝对收敛。

Fourier 级数要解决的问题是，对于给定的 $2\pi$ 周期函数 $f(x)$，三角级数
$$
f(x) = \frac{a_0}{2} + \sum_{n = 1}^{\infty} (a_n \cos n x + b_n \sin n x)
$$
是否存在？系数取值是什么？

我们假设存在这么一个级数一致收敛于 $f(x)$，那么：

\begin{enumerate}
    \item 求 $a_0$：
    $$
    \begin{aligned}
        \int_{-\pi}^{\pi} f(x) \dd{x} & = \int_{-\pi}^{\pi} \left( \frac{a_0}{2} + \sum_{n = 1}^{\infty} (a_n \cos n x + b_n \sin n x) \right) \dd{x} \\
        & = \int_{-\pi}^{\pi} \frac{a_0}{2} \dd{x} + \sum_{n = 1}^{\infty} \int_{-\pi}^{\pi} (a_n \cos n x + b_n \sin n x) \dd{x} \\
        & = \pi a_0
    \end{aligned}
    $$
    所以
    $$
    a_0 = \frac{1}{\pi} \int_{-\pi}^{\pi} f(x) \dd{x}
    $$
    \item 求 $a_n$：
    $$
    \begin{aligned}
        \int_{-\pi}^{\pi} f(x) \cos n x \dd{x} & = \int_{-\pi}^{\pi} \ab(\frac{a_0}{2} \cos n x + \sum_{k = 1}^{\infty} (a_k \cos k x \cos n x + b_k \sin k x \cos n x)) \dd{x} \\
        & = \frac{a_0}{2} \int_{-\pi}^{\pi} \cos n x \dd{x} + \sum_{k = 1}^{\infty} \int_{-\pi}^{\pi} (a_k \cos k x \cos n x + b_k \sin k x \cos n x) \dd{x} \\
        & = \pi a_n
    \end{aligned}
    $$
    所以
    $$
    a_n = \frac{1}{\pi} \int_{-\pi}^{\pi} f(x) \cos n x \dd{x}
    $$
    \item 求 $b_n$：
    $$
    \begin{aligned}
        \int_{-\pi}^{\pi} f(x) \sin n x \dd{x} & = \int_{-\pi}^{\pi} \ab(\frac{a_0}{2} \sin n x + \sum_{k = 1}^{\infty} (a_k \cos k x \sin n x + b_k \sin k x \sin n x)\ab) \dd{x} \\
        & = \frac{a_0}{2} \int_{-\pi}^{\pi} \sin n x \dd{x} + \sum_{k = 1}^{\infty} \int_{-\pi}^{\pi} (a_k \cos k x \sin n x + b_k \sin k x \sin n x) \dd{x} \\
        & = \pi b_n
    \end{aligned}
    $$
    所以
    $$
    b_n = \frac{1}{\pi} \int_{-\pi}^{\pi} f(x) \sin n x \dd{x}
    $$
\end{enumerate}

\begin{definition}[Fourier 级数]
    对于周期为 $2\pi$ 的周期函数 $f(x)$，定义其 \textbf{Fourier 级数}为：
    $$
    f(x) = \frac{a_0}{2} + \sum_{n = 1}^{\infty} (a_n \cos n x + b_n \sin n x)
    $$
    其中：
    $$
    \begin{gathered}
        a_n = \frac{1}{\pi} \int_{-\pi}^{\pi} f(x) \cos n x\ (n \in \mathbb{N}) \\
        b_n = \frac{1}{\pi} \int_{-\pi}^{\pi} f(x) \sin n x\ (n \in \mathbb{N}^+)
    \end{gathered}
    $$
\end{definition}

Fourier 级数的定义仅要求 $f(x)$ 可积，但是级数的取值和函数的光滑性有关。

\begin{definition}[光滑]
    称 $f(x)$ 在 $[l, r]$ \textbf{光滑}当且仅当 $f'(x)$ 在 $[l, r]$ 上连续。称 $f(x)$ 在 $[l, r]$ \textbf{分段光滑}当且仅当 $f'(x)$ 在 $[l, r]$ 仅在有限个第一类间断点处不连续。
\end{definition}

\begin{theorem}[Dirichlet 定理]
    若 $f(x)$ 是以 $2\pi$ 为周期的周期函数，且在 $[-\pi, \pi]$ 分段光滑，那么对于 $x \in [-\pi, \pi]$，$f(x)$ 的 Fourier 级数收敛于 $f(x)$ 左右极限的平均值。
\end{theorem}